% Um pregão passo a passo, com os números reais do robô (configuração ARIMA-GARCH
% Sharpe, traço da refação de 05/10/2026). Dois pregões seguidos com o mesmo
% ARIMA(3,1,0): no primeiro o GARCH libera a compra, no segundo veta.
\documentclass[border=6pt]{standalone}
% Preâmbulo comum dos gráficos que explicam os métodos (30/10/2026).
% Dados e números: dados/resultados/metodos_ilustrados_2026-10-30/, gerados por
% ferramentas/MetodosIlustrados.mq5 com as funções do próprio robô. Nenhum
% número é digitado aqui: os do texto vêm de valores.tex (\vX para o texto,
% \pX para coordenada).
\usepackage[T1]{fontenc}
\usepackage{amsmath}
\usepackage{mathptmx}  % Times, como as figuras do TG
\usepackage{pgfplots}
\usepgfplotslibrary{groupplots,dateplot,colormaps,fillbetween}
\usetikzlibrary{arrows.meta,positioning}
\pgfplotsset{compat=1.18}
\definecolor{cmod}{RGB}{0,90,181}     % modelo / ajuste
\definecolor{cdes}{RGB}{220,50,32}    % destaque, limite, valor crítico
\definecolor{cdad}{RGB}{120,120,120}  % dado observado
\newcommand{\arq}[1]{../../dados/resultados/metodos_ilustrados_2026-10-30/#1}
\input{\arq{valores.tex}}
% Figuras 08 a 14 (o robô operando): traço por barra da refação de 05/10/2026,
% resumido por ferramentas/scripts/robo_em_acao.sh. Macros \vR e \pR.
\newcommand{\arqr}[1]{../../dados/resultados/robo_em_acao_2026-10-30/#1}
\input{\arqr{valores.tex}}
\definecolor{cari}{RGB}{230,140,0}   % ARIMA puro
\definecolor{cok}{RGB}{0,140,70}     % compra liberada
\usetikzlibrary{shapes.misc}
\newcommand{\passo}[1]{\tikz[baseline=(n.base)]\node[circle,fill=black!75,text=white,inner sep=1pt,font=\footnotesize\bfseries](n){#1};}
% eixo de datas comum às figuras do traço (jul/2019 a jun/2021)
\pgfplotsset{traco/.style={date coordinates in=x, xmin=2019-07-01, xmax=2021-06-30,
    xtick={2019-07-01,2020-01-01,2020-07-01,2021-01-01,2021-06-30},
    xticklabels={jul/2019,jan/2020,jul/2020,jan/2021,jun/2021}}}
\pgfplotsset{
  base/.style={table/col sep=comma, grid=major, grid style={gray!20},
    tick label style={font=\small}, label style={font=\small},
    title style={font=\small\bfseries, align=center},
    every axis plot/.append style={line width=0.9pt},
    /pgf/number format/use comma, /pgf/number format/1000 sep={.}},
}
% faixa vertical nos pregões do traço em que a condição #1 vale, na cor #2,
% do piso #3 ao teto #4 do eixo
\newcommand{\faixa}[4]{%
  \addplot[name path=topo, draw=none, const plot, forget plot] table[x=data, y expr={(#1) ? #4 : #3}]{\arqr{dia_a_dia.csv}};
  \addplot[name path=piso, draw=none, forget plot] coordinates {(2019-07-01,#3) (2021-06-30,#3)};
  \addplot[#2, forget plot] fill between[of=topo and piso];}

\begin{document}
\begin{tikzpicture}[
  rotulo/.style={anchor=west, text width=4.3cm, font=\small, align=left},
  caixa/.style={draw=#1, fill=#1!6, rounded corners=3pt, text width=6.1cm, minimum height=1.25cm,
    font=\small, align=left, inner sep=5pt},
  caixa/.default=cmod, seta/.style={-{Stealth[length=5pt]}, black!50, line width=0.8pt}]
\def\dy{1.75}
% cabeçalhos
\node[font=\bfseries, align=center] at (7.6,1.35) {\vRDiaLib\\[-1pt]{\normalfont\small (sexta-feira)}};
\node[font=\bfseries, align=center] at (14.4,1.35) {\vRDiaVet\\[-1pt]{\normalfont\small (o pregão seguinte)}};
% passos
\foreach \i/\txt in {
  1/{\passo{1} \textbf{Dados}\\240 fechamentos, até a véspera},
  2/{\passo{2} \textbf{Teste ADF}\\(refeito a cada 21 pregões)},
  3/{\passo{3} \textbf{Grade do AIC}\\(refeita a cada 21 pregões)},
  4/{\passo{4} \textbf{Hannan-Rissanen}\\(todo pregão)},
  5/{\passo{5} \textbf{Previsão do ARIMA}\\$\Rightarrow$ direção},
  6/{\passo{6} \textbf{GARCH(1,1)} nos resíduos\\$\Rightarrow$ risco},
  7/{\passo{7} \textbf{Decisão}},
  8/{\passo{8} \textbf{O que aconteceu}\\(não é previsão do robô)}}
  \node[rotulo] (r\i) at (0,-\i*\dy+\dy) {\txt};
% coluna do pregão liberado
\node[caixa=cdad] (a1) at (7.6,0) {Último fechamento: R\$ \vRPrecoLib\\Série usada: log do preço};
\node[caixa=cdad] (a2) at (7.6,-\dy) {Log-preço não estacionário $\Rightarrow$ $d = \vRDLib$\\O modelo trabalha com o log-retorno};
\node[caixa=cdad] (a3) at (7.6,-2*\dy) {Menor AIC: ARIMA(\vRPLib,\vRDLib,\vRQLib)\\três retornos passados, nenhum erro};
\node[caixa=cmod] (a4) at (7.6,-3*\dy) {$\phi_1 = \vRPhiUmLib$, $\phi_2 = \vRPhiDoisLib$, $\phi_3 = \vRPhiTresLib$\\constante $= \vRConstLib\%$ ao dia};
\node[caixa=cmod] (a5) at (7.6,-4*\dy) {Retorno previsto: $+\vRPrevLib\%$\\Positivo $\Rightarrow$ sinal de \textbf{compra}};
\node[caixa=cmod] (a6) at (7.6,-5*\dy) {$\alpha = \vRAlfaLib$, $\beta = \vRBetaLib$\\$\sigma$ prevista $\vRVolLib\%$ $\le$ limiar $\vRLimLib\%$};
\node[caixa=cok] (a7) at (7.6,-6*\dy) {\textbf{Liberada}: o robô compra na abertura};
\node[caixa=cdad] (a8) at (7.6,-7*\dy) {Variação do dia: $\vRRealLib\%$\\(o robô está comprado)};
% coluna do pregão vetado
\node[caixa=cdad] (b1) at (14.4,0) {Último fechamento: R\$ \vRPrecoVet\\(a queda de sexta já entra na janela)};
\node[caixa=cdad] (b2) at (14.4,-\dy) {Mesmo $d = \vRDVet$};
\node[caixa=cdad] (b3) at (14.4,-2*\dy) {Mesmo ARIMA(\vRPVet,\vRDVet,\vRQVet)};
\node[caixa=cmod] (b4) at (14.4,-3*\dy) {$\phi_1 = \vRPhiUmVet$, $\phi_2 = \vRPhiDoisVet$, $\phi_3 = \vRPhiTresVet$\\coeficientes reestimados, quase iguais};
\node[caixa=cmod] (b5) at (14.4,-4*\dy) {Retorno previsto: $+\vRPrevVet\%$\\Positivo $\Rightarrow$ sinal de \textbf{compra}};
\node[caixa=cdes] (b6) at (14.4,-5*\dy) {$\alpha = \vRAlfaVet$, $\beta = \vRBetaVet$\\$\sigma$ prevista $\vRVolVet\%$ $>$ limiar $\vRLimVet\%$};
\node[caixa=cdes] (b7) at (14.4,-6*\dy) {\textbf{Vetada}: não compra, e a posição de\\sexta é zerada na abertura};
\node[caixa=cdad] (b8) at (14.4,-7*\dy) {Variação do dia: $\vRRealVet\%$. O ARIMA puro,\\sem o veto, só saiu em \vRDiaSaidaA};
\foreach \i [evaluate=\i as \j using int(\i+1)] in {1,...,7} {
  \draw[seta] (a\i) -- (a\j); \draw[seta] (b\i) -- (b\j); }
% nota
\node[font=\footnotesize, text=black!75, align=left, anchor=north west, text width=17.6cm]
  at (0,-7*\dy-0.95) {Os dois pregões usam o mesmo modelo e preveem alta. A diferença está só no passo 6: a queda de sexta
  entrou nos resíduos, o GARCH elevou a volatilidade prevista acima do limiar (percentil 70 da volatilidade ajustada
  na janela) e vetou. O GARCH reage ao que já aconteceu: ele não previu a queda de segunda, e a sexta foi perdida do mesmo jeito.
  Este é um exemplo escolhido para mostrar o mecanismo, não um resultado médio.};
\end{tikzpicture}
\end{document}
